\section*{Beitrag 5 - Prof. Dr. David Garcia - 
Understanding Emotions through Digital Traces}

The main topic is how digital traces left from people in social media platforms is used to study collective emotions using computational social science.
This includes specific events such as collective emotions based on twitter posts and YouGov surveys like a terrorist attack in France, the George Floyd Incident and many more.

The talk further handles the topic of how these posts are actually used to extract insights. We are shown how the tools and methods are validated by comparing them to established social science methods such as real life surveys.

The extraction and storage of such large amounts of data is related with big data management, while network visualization is used to extract exploratory insights about
collective emotions. 
Machine learning models like LIWC or more advanced models like LEIA are also discussed as a way of labeling the underlying data, which is infeasible to do manually due to the large amount of data.
The inherent digital nature of the collected data allows scientists to produce results based on an amount of data that would have been difficult to analyze in the past. 
Modern data collection makes access to data easier and provides higher temporal resolution compared to real study settings.

In my computer science studies I came in contact with social media networks in a network visualizations course, making it particurly interesting for me to see the methods I learned used in a research setting.
Additionally, how a model is trained with training datasets and then has is generizability tested is also an aspect I came in contact with my data mining and Introduction to Machine Learning courses.

Understanding how a population feels about a certain event, or predicting a certain response to an event, can offer goverments a good way to handle local as well as global crises. 
This field of research might be in particular interest to me, because I am contemplating turning my focus to social studies. 

On a personal level I was also motivated to think about what kind of digital traces I leave when I use the internet during the lecture.